\documentclass[12pt]{article}

\usepackage[a4paper, margin=2.5cm]{geometry}
\usepackage{amsmath}
\usepackage{amssymb}
\usepackage{amsthm}
\usepackage[colorlinks=true, linkcolor=blue, citecolor=blue, urlcolor=blue]{hyperref}
\usepackage{booktabs}
\usepackage{natbib}

\newtheorem{theorem}{Theorem}
\newtheorem{proposition}[theorem]{Proposition}
\newtheorem{definition}[theorem]{Definition}
\newtheorem{conjecture}[theorem]{Conjecture}
\newtheorem{remark}[theorem]{Remark}

\newcommand{\Mcal}{\mathcal{M}}
\newcommand{\phig}{\varphi}
\newcommand{\ds}{d_s}
\newcommand{\dseff}{d_s^{\mathrm{eff}}}

\title{Spectral Structure of the 600-Cell Vertex Graph:\\Eigenvalue Connections and the F4 Fibonacci Extension}
\author{%
  Lee Smart\\[4pt]
  \textit{Institute of Vibrational Field Dynamics}\\[2pt]
  \texttt{contact@vibrationalfielddynamics.org}\\[2pt]
  \texttt{@vfd\_org}%
}
\date{April 2026}

\begin{document}

\maketitle

\noindent\textbf{Status:} Preprint (not peer-reviewed)

\begin{abstract}
We report two spectral results connecting the 600-cell to the closure-class mass programme. First, the graph Laplacian of the 600-cell vertex adjacency graph ($|V| = 120$, degree 12) has eigenvalue $\lambda = 15$ with multiplicity 16. This equals $\Delta C = 15$, the closure invariant difference between proton and electron, embedding the mass law's leading-order exponent directly in the polytope spectrum. This is a theorem-level connection, verifiable by diagonalising the $120 \times 120$ Laplacian. Second, the F4 Fibonacci hierarchical extension of the $\phig$-manifold --- a construction in which shell $k$ contains $F(k{+}2)$ sub-cells --- produces an effective spectral dimension converging toward $\phig$ from above: $\dseff^{HK} = 1.635$ at $V = 4960$, with empirical $1/V$ extrapolation giving $\ds(\infty) = 1.628$ ($R^2 = 0.996$). Among all tested control constructions, F4 is the only one exhibiting this convergence. An earlier claim that the 600-cell vertex graph itself has effective spectral dimension $\dseff \approx 1.637$--$1.642$ near $\phig$ has been retracted: finer-resolution analysis shows this was an artifact of coarse time-windowing across mixed spectral regimes.
\end{abstract}

%==================================================================
\section{Introduction}\label{sec:intro}
%==================================================================

This paper presents two spectral results relevant to the closure-class mass programme~\citep{SmartI, SmartII}. The programme requires a spectral connection between the 600-cell geometry and the golden ratio $\phig$. An earlier version of this paper reported that the effective spectral dimension of the 600-cell vertex graph lies near $\phig$; that claim is retracted here with explanation (Section~\ref{sec:retraction}).

What survives and what is new:

\begin{enumerate}
    \item \textbf{Eigenvalue connection (theorem).} The 600-cell Laplacian has eigenvalue $\lambda = 15$ with multiplicity 16. This equals $\Delta C = 15$, the closure invariant that governs the leading-order proton-to-electron mass ratio. The mass exponent is not merely a combinatorial quantity: it is a specific point in the polytope spectrum. (Section~\ref{sec:eigenvalue}.)

    \item \textbf{F4 spectral dimension (numerical).} The F4 Fibonacci hierarchical construction produces $\dseff \to \phig$ as graph size increases, with $\dseff = 1.635$ at $V = 4960$. This is independent of the 600-cell spectral-dimension measurement and provides the primary numerical evidence for the conjectural identification $\ds = \phig$ used in the lepton winding operator. (Section~\ref{sec:F4}.)

    \item \textbf{Retraction of the 600-cell $\dseff \approx \phig$ claim.} The earlier measurement was an artifact of coarse time-windowing. Finer analysis shows $\dseff(t)$ crosses $\phig$ at specific times but does not plateau there. (Section~\ref{sec:retraction}.)
\end{enumerate}

%==================================================================
\section{The 600-Cell Laplacian Eigenvalue $\lambda = 15$}\label{sec:eigenvalue}
%==================================================================

\subsection{The 600-cell vertex graph}

The 600-cell is the regular $H_4$ polytope with 120 vertices, 720 edges, 1200 triangular faces, and 600 tetrahedral cells~\citep{Coxeter1973}. Its vertex adjacency graph $G_V$ has $|V| = 120$ vertices, each of degree 12. The vertex coordinates are expressible in terms of $\phig$.

\subsection{Eigenvalue spectrum}

The graph Laplacian $L = D - A$ of $G_V$ has distinct eigenvalues. Direct diagonalisation yields the eigenvalue $\lambda = 15$ with multiplicity 16.

\begin{proposition}[Spectral embedding of $\Delta C$]\label{prop:eigenvalue}
The graph Laplacian of the 600-cell vertex adjacency graph has eigenvalue $\lambda = 15$. The closure invariant difference between proton and electron closure classes is $\Delta C = C(\{2,3,4\}) - C(\{1\}) = 14 - (-1) = 15$. Therefore $\Delta C$ is an eigenvalue of the 600-cell Laplacian.
\end{proposition}

\textsc{Status: theorem.} Both the eigenvalue and the invariant difference are verifiable by direct computation. The connection requires no statistical estimation, windowing, or modelling.

\subsection{Significance}

The mass law's leading-order term $\phig^{\Delta C} = \phig^{15}$ is not merely a combinatorial formula: the exponent $15$ is a specific eigenvalue of the Laplacian of the geometry that defines the framework. This provides a direct spectral link between the mass law and the polytope, independent of the effective spectral dimension.

\subsection{Additional $\phig$-connections}

The 600-cell is connected to $\phig$ through multiple channels:
\begin{itemize}
    \item Vertex coordinates expressible in terms of $\phig$~\citep{Coxeter1973}.
    \item $E_8$ Coxeter projection: 240 roots of $E_8$ project onto 120 vertices of the 600-cell.
    \item Spectral gap of the 5-shell path graph: $\lambda_1(P_5) = 2 - \phig = \phig^{-2}$ (algebraic identity).
    \item Five distance classes in the vertex graph ($N = 5$), matching the shell count.
\end{itemize}

%==================================================================
\section{F4 Fibonacci Hierarchical Extension}\label{sec:F4}
%==================================================================

\subsection{Construction}

The F4 graph $G_F(N, b)$ assigns $F(k{+}2)$ sub-cells to shell $k$ (where $F$ is the Fibonacci sequence), each sub-cell containing $b$ vertices. Intra-cell edges form complete subgraphs ($K_b$, weight 1). Intra-shell ring edges connect sub-cells cyclically (weight $1/\phig$). Inter-shell bridge edges connect adjacent shells (weight 1). For $N = 5$: total $|V| = 31b$.

\textsc{Status:} The ring weight $1/\phig$ and the shell count $N = 5$ are construction inputs, not derived from the base formalism.

\subsection{Spectral dimension convergence}

\begin{table}[ht]
\centering
\caption{Effective spectral dimension of the F4 family.}
\label{tab:F4}
\begin{tabular}{@{}rrcrc@{}}
\toprule
$b$ & $V$ & $\dseff^{HK}$ & $\dseff^{RW}$ & HK--RW gap \\
\midrule
6 & 186 & 1.867 & 1.834 & 0.034 \\
10 & 310 & 1.787 & 1.780 & 0.008 \\
30 & 930 & 1.679 & 1.684 & 0.005 \\
80 & 2480 & 1.646 & 1.645 & 0.001 \\
160 & 4960 & 1.635 & --- & --- \\
\bottomrule
\end{tabular}
\end{table}

Empirical $1/V$ extrapolation: $\ds(\infty) = 1.628$ with $R^2 = 0.996$. The distance $|\ds(\infty) - \phig| = 0.010$.

\subsection{Control constructions}

\begin{table}[ht]
\centering
\caption{Control constructions: none trend toward $\phig$.}
\label{tab:controls}
\begin{tabular}{@{}lccc@{}}
\toprule
Construction & $V$ & $\dseff^{RW}$ & Approaches $\phig$? \\
\midrule
F4 Fibonacci & $\sim\!5000$ & $\sim\!1.64$ & Yes (descending) \\
Binary tree & 2047 & 1.806 & No ($\to$ 2) \\
Ternary tree & 1093 & 1.783 & No ($\to$ 2) \\
Non-Fibonacci hierarchical & 480 & 2.038 & No ($\to$ 2) \\
1D chain & 250 & 0.919 & No ($\to$ 1) \\
\bottomrule
\end{tabular}
\end{table}

Among the tested control families, F4 is the only one whose effective spectral dimension trends toward $\phig$.

\subsection{Perturbation robustness}

The random-walk spectral dimension estimate is robust to $20\%$ edge-weight perturbation ($\Delta \dseff \approx 0.001$ at $V = 2480$).

\subsection{Conjecture}

\begin{conjecture}[F4 Spectral Dimension]\label{conj:F4}
$\dseff(G_F(5,b)) \to \phig$ as $b \to \infty$.
\end{conjecture}

\textsc{Status: conjectured.} Supported by monotone convergence from above, HK--RW agreement, control separation, and perturbation robustness. Not proven analytically.

%==================================================================
\section{Retraction of the 600-Cell $\dseff \approx \phig$ Claim}\label{sec:retraction}
%==================================================================

An earlier version of this paper reported that the effective spectral dimension of the 600-cell vertex graph lies in the interval $\dseff \approx 1.637$--$1.642$, near $\phig \approx 1.618$. This claim is retracted.

\subsection{What went wrong}

A more careful analysis with finer time resolution (10,000 time points vs.\ the original coarser grid) reveals that $\dseff(t)$ on the 600-cell vertex graph does not plateau near $\phig$. The function $\dseff(t)$ crosses $\phig$ at two specific times ($t \approx 0.074$ and $t \approx 0.98$) but the stable plateau is at $\dseff \approx 3.6$, consistent with $2 + \phig \approx 3.618$. The earlier reported value of $\approx 1.637$ was the median of a mixed distribution across spectral regimes --- an artifact of the time-windowing methodology, not a genuine plateau.

\subsection{What survives}

\begin{itemize}
    \item The F4 construction (Section~\ref{sec:F4}) is entirely independent of the 600-cell spectral dimension measurement. Its convergence toward $\phig$ is established by proper plateau detection at each graph size.
    \item The eigenvalue connection $\lambda = 15 = \Delta C$ (Section~\ref{sec:eigenvalue}) is a theorem, unaffected by the spectral dimension retraction.
    \item The spectral gap $\lambda_1(P_5) = \phig^{-2}$ on the 5-shell path graph is an algebraic identity, unaffected.
\end{itemize}

\subsection{Lessons}

The retraction illustrates the importance of fine-grained plateau detection in heat-kernel spectral dimension estimation on finite graphs. Coarse time-windowing can produce apparent plateaux that are artifacts of mixing across spectral regimes. The F4 analysis avoids this by using scale-dependent plateau detection across a family of increasing graph sizes.

%==================================================================
\section{Bridge to the Mass Programme}\label{sec:bridge}
%==================================================================

The programme linkage after the retraction:
\begin{center}
\begin{tabular}{@{}ll@{}}
Paper~I & proton/electron mass law (derived within framework) \\
Paper~II & no-go + conditional winding bridge \\
Paper~III (this paper) & eigenvalue $\lambda = 15 = \Delta C$ (theorem) \\
 & F4 spectral dimension $\to \phig$ (numerical) \\
 & 600-cell $\dseff \approx \phig$ (retracted) \\
\end{tabular}
\end{center}

The spectral identification $\ds = \phig$ required by Paper~II's winding operator now rests on the F4 evidence alone, not on the 600-cell. The 600-cell contributes to the mass programme through the eigenvalue connection and through the structural presence of $\phig$ in its geometry, not through its effective spectral dimension.

%==================================================================
\section{Claim Status}\label{sec:status}
%==================================================================

\begin{table}[ht]
\centering
\caption{Claim status table.}
\label{tab:status}
\begin{tabular}{@{}lll@{}}
\toprule
Claim & Status & Basis \\
\midrule
$\lambda = 15$ is a 600-cell eigenvalue & \textsc{theorem} & Laplacian diagonalisation \\
$\Delta C = 15$ & \textsc{theorem} & Arithmetic from closure invariant \\
$\lambda_1(P_5) = \phig^{-2}$ & \textsc{theorem} & $2 - 2\cos(\pi/5) = 2 - \phig$ \\
F4: $\dseff \to \phig$ & \textsc{conjectured} & Numerical convergence ($R^2 = 0.996$) \\
F4: perturbation robust & \textsc{numerical observation} & $\Delta \dseff < 0.002$ at $20\%$ \\
600-cell $\dseff \approx \phig$ & \textsc{retracted} & Windowing artifact \\
$\ds = \phig$ (exact) & \textsc{open} & No analytic proof \\
\bottomrule
\end{tabular}
\end{table}

%==================================================================
\section{Open Problems}\label{sec:open}
%==================================================================

\begin{enumerate}
    \item \textbf{Analytic proof of Conjecture~\ref{conj:F4}.} Does $\dseff(G_F(5,b)) \to \phig$ as $b \to \infty$?
    \item \textbf{Structural origin of $\lambda = 15$.} Why does the 600-cell Laplacian contain $\Delta C$ as an eigenvalue? Is this forced by the $H_4$ representation theory?
    \item \textbf{Connection between eigenvalue and spectral dimension.} The eigenvalue $\lambda = 15 = \Delta C$ and the F4 spectral dimension convergence toward $\phig$ are currently independent observations. Whether they share a common spectral mechanism is unknown.
    \item \textbf{Derivation of $\phig^5$.} The coefficient $\phig^5$ in the winding operator is structurally selected but not derived.
\end{enumerate}

%==================================================================
\section{Conclusion}\label{sec:conclusion}
%==================================================================

The 600-cell connects to the mass programme through a theorem-level spectral link: its Laplacian eigenvalue $\lambda = 15$ equals the closure invariant difference $\Delta C = 15$ that governs the leading-order proton-to-electron mass ratio. This connection is exact and requires no numerical estimation.

The F4 Fibonacci hierarchical extension provides independent numerical evidence that a $\phig$-structured geometry produces $\dseff \to \phig$, supporting the conjectural spectral identification used in the lepton winding operator. An earlier claim that the 600-cell vertex graph itself has $\dseff \approx \phig$ has been retracted: the reported value was an artifact of coarse time-windowing. The retraction does not affect the eigenvalue theorem, the F4 convergence, or the spectral gap identity $\lambda_1(P_5) = \phig^{-2}$.

%==================================================================
\bibliographystyle{plainnat}
\bibliography{references}

\end{document}
